% longtable cannot break pages in two-column mode, which is what acmart
% sigconf sets. The appendix therefore runs single-column, which also suits
% six-column tables better than a 3.33in column would.
\clearpage
\onecolumn

\section{The full census}
\label{sec:census}

Section~\ref{sec:taxonomy} describes the largest families and counts the rest.
This appendix is the rest: one row for every named mechanism and every named
defense the corpus contains, with the coding each was given. It is here so that
the claim behind the counts --- that both literatures were read in full rather
than sampled where convenient --- can be checked rather than taken on trust.

Three conventions are worth stating before the tables.

\textbf{Entries whose authors named nothing are included, in italics.} The
tables in the body exclude them, because an unnamed row makes a poor example.
A census is the opposite case: omitting them would undercount the corpus and
overstate how much of the field names its own contributions. That share is
itself a finding, and it differs between the two sides.

\textbf{\emph{Validated against} is what each paper's own evaluation tested},
read from its evaluation section, not what its abstract claims. Where the two
disagreed we coded the evaluation. This is the field that carries
Section~\ref{sec:emptycell}, so it was the one we checked hardest.

\textbf{Every row cites its source.} All 715 distinct papers behind these
entries resolve to a bibliography key, so no named contribution appears here
without attribution. The machine-readable registries released with this work
carry more than a citation can: each entry's DOI or arXiv identifier, the
quoted evidence its coding was taken from, and the coder's confidence.

% GENERATED by scripts/gen_census_appendix.py -- do not hand-edit.
\begingroup
\footnotesize
\setlength{\tabcolsep}{3.5pt}
\begin{longtable}{p{0.33\textwidth}llll}
\caption{Every named mechanism in the corpus (233 entries). Rows in \emph{italics} are entries whose authors gave their contribution no name; they are counted here because excluding them would overstate how much of the field names what it finds. Source is first author and year; full records with DOI or arXiv identifiers ship with the artifact.}\
\label{tab:census-mech}\\
\toprule
\textbf{Mechanism or attack} & \textbf{Channel} & \textbf{Consequence} & \textbf{Track} & \textbf{Source} \\
\midrule
\endfirsthead
\multicolumn{5}{l}{\emph{Table~\ref{tab:census-mech}: named mechanisms, continued}} \\
\toprule
\textbf{Mechanism or attack} & \textbf{Channel} & \textbf{Consequence} & \textbf{Track} & \textbf{Source} \\
\midrule
\endhead
\midrule
\multicolumn{5}{r}{\emph{continued}} \\
\endfoot
\bottomrule
\endlastfoot
lost in the middle & direct-input & reasoning-corruption & ML/AI & Liu 2023 \\
distraction by irrelevant context & direct-input & reasoning-corruption & ML/AI & Shi 2023 \\
\emph{context-length-induced degradation independent of retrieval failure} & direct-input & reasoning-corruption & ML/AI & Du 2025 \\
context drift & direct-input & reasoning-corruption & ML/AI & Dongre 2025 \\
\emph{getting lost in multi-turn conversation (premature-answer unreliability under underspecification)} & direct-input & reasoning-corruption & ML/AI & Laban 2025 \\
classifier context rot & direct-input & reasoning-corruption & ML/AI & Martin 2026 \\
mechanistic explanation of lost-in-the-middle (training-induced primacy/recency from retrieval-demand adaptation) & direct-input & reasoning-corruption & ML/AI & Salvatore 2025 \\
Attention Latch & direct-input & reasoning-corruption & ML/AI & Shehata 2026 \\
mechanistic refinement of lost-in-the-middle (threshold-dependent onset at \textasciitilde{}50\% of context window; effect vanishes near window limits) & direct-input & reasoning-corruption & ML/AI & Veseli 2025 \\
mechanistic explanation of lost-in-the-middle (positional hidden-state channels, not just position embeddings) & direct-input & reasoning-corruption & ML/AI & Yu 2024 \\
\emph{long-context safety/refusal degradation in agents} & direct-input & reasoning-corruption & ML/AI & Hadeliya 2025 \\
cognitive bandwidth bottleneck (action-space scaling) & direct-input & reasoning-corruption & ML/AI & Xu 2025 \\
\emph{abrupt persona/tone instability from lack of affective state inertia} & direct-input & reasoning-corruption & ML/AI & Subaharan 2026 \\
Governance Decay & direct-input & silent-corruption & ML/AI & Chen 2026 \\
\emph{unverified success fabrication in unattended long-horizon agents (agent honesty gap)} & direct-input & silent-corruption & ML/AI & Deng 2026 \\
Context Pollution (evolutionary search history bias) & direct-input & reasoning-corruption & ML/AI & Yan 2026 \\
\emph{plan non-persistence under context eviction (context-time memory volatility)} & direct-input & silent-corruption & ML/AI & Mehta 2026 \\
context proprioceptive blindness & direct-input & reasoning-corruption & ML/AI & Xu 2026 \\
Distributional Collapse & direct-input & silent-corruption & ML/AI & Ebrahimzadeh 2026 \\
\emph{summary-length invariance under context compaction} & direct-input & reasoning-corruption & ML/AI & Cim 2026 \\
Mandate Salience Decay (MSD) & direct-input & reasoning-corruption & ML/AI & Safder 2026 \\
Contextual Belief Management (CBM) failure & direct-input & reasoning-corruption & ML/AI & Xu 2026 \\
agent drift (overthinking / overacting) & direct-input & reasoning-corruption & ML/AI & Sui 2026 \\
context pollution (role conflation in single-agent code-as-action agents) & direct-input & silent-corruption & ML/AI & Fei 2026 \\
cumulative contextual decay (attention pollution, dilution, and drift) & direct-input & reasoning-corruption & ML/AI & Hong 2025 \\
contextual misguidance & direct-input & silent-corruption & ML/AI & Wang 2026 \\
superposition of contextual information and parametric memory (mechanistic explanation of knowledge conflict) & direct-input & silent-corruption & ML/AI & Li 2025 \\
position exhaustion & direct-input & reasoning-corruption & ML/AI & Huo 2026 \\
\emph{"malloc-only" context accumulation causing reasoning degeneration in extended CoT} & direct-input & reasoning-corruption & ML/AI & Zheng 2026 \\
scaffolding collapse & direct-input & reasoning-corruption & ML/AI & Shao 2026 \\
KV-cache contamination & direct-input & reasoning-corruption & ML/AI & Kang 2026 \\
Activation-guided MCMC direct prompt injection (EBM-guided token-level MCMC; no acronym given) & direct-input & goal-hijack & Security & Li 2025 \\
ASTRA (and ASTRA++) & direct-input & goal-hijack & Security & Pandya 2025 \\
Blind attacks & direct-input & goal-hijack & Security & Liu 2025 \\
Brain-Prompt Injection & direct-input & goal-hijack & Security & Tai 2026 \\
Cognitive Overload Attack & direct-input & goal-hijack & Security & Upadhayay 2024 \\
Combined Attack (Open-Prompt-Injection) & direct-input & goal-hijack & Security & Liu 2023 \\
CoT Forgery (Chain-of-Thought Forgery) & direct-input & goal-hijack & Security & Ye 2026 \\
CyberSecEval 2: prompt-injection + code-interpreter-abuse test set & direct-input & goal-hijack & Security & Bhatt 2024 \\
distraction effect (attention shift to injected instruction) & direct-input & reasoning-corruption & Security & Hung 2024 \\
G2PIA (goal-guided generative prompt injection) & direct-input & goal-hijack & Security & Zhang 2024 \\
Guidance Injection & direct-input & persistence-backdoor & Security & Liu 2026 \\
HackAPrompt: prompt-hacking taxonomy and 600k-prompt corpus & direct-input & goal-hijack & Security & Schulhoff 2023 \\
HOUYI & direct-input & goal-hijack & Security & Liu 2023 \\
JudgeDeceiver & direct-input & goal-hijack & Security & Shi 2024 \\
LLMSmith: prompt-injection-to-RCE exploitation & direct-input & goal-hijack & Security & Liu 2023 \\
message history forging (direct prompt injection via forged system/assistant roles) & direct-input & goal-hijack & Security & Kaya 2025 \\
Overflow & direct-input & resource-abuse & Security & Feiglin 2026 \\
P2SQL injection (prompt-to-SQL) & direct-input & goal-hijack & Security & Pedro 2023 \\
POUGH & direct-input & goal-hijack & Security & Huang 2024 \\
preemptive answer attack (CoT short-circuiting) & direct-input & reasoning-corruption & Security & Xu 2024 \\
Reasoning Interruption Attack (adaptive token compression) & direct-input & reasoning-corruption & Security & Cui 2025 \\
reasoning token overflow (RTO) / practical reasoning interruption attack & direct-input & reasoning-corruption & Security & Cui 2025 \\
Reverse CAPTCHA (invisible Unicode instruction injection) & direct-input & goal-hijack & Security & Graves 2026 \\
Separator Injection Attack & direct-input & goal-hijack & Security & Li 2025 \\
situation-driven contextual adversarial prompt attack & direct-input & goal-hijack & Security & Das 2024 \\
structured attack generation framework & direct-input & goal-hijack & Security & Cai 2026 \\
System Prompt Poisoning (SPP) & direct-input & persistence-backdoor & Security & Guo 2025 \\
Tensor Trust: prompt extraction / prompt hijacking & direct-input & goal-hijack & Security & Toyer 2023 \\
\emph{adaptive attack suite bypassing 12 PI/jailbreak defenses} & direct-input & goal-hijack & Security & Nasr 2025 \\
\emph{adaptive prompt-injection attacks for defense evaluation (PIEval)} & direct-input & goal-hijack & Security & Jia 2025 \\
\emph{automatic universal gradient-based prompt injection} & direct-input & goal-hijack & Security & Liu 2024 \\
\emph{diverse RL red-teaming (auto-generated rewards + multi-step RL)} & direct-input & goal-hijack & Security & Beutel 2024 \\
\emph{Prompt Injection 2.0 (hybrid PI + classic web-exploit combination attacks)} & direct-input & goal-hijack & Security & McHugh 2025 \\
\emph{rule-tag semantic reframing template attack ('Breaking the Prompt Wall')} & direct-input & goal-hijack & Security & Chang 2025 \\
\emph{system-prompt and uploaded-file extraction from custom GPTs} & direct-input & data-exfiltration & Security & Yu 2023 \\
Xmera (chimera MitM framework) & direct-input & silent-corruption & Security & Fastowski 2025 \\
context rot (agentic/long-horizon search) & tool-output & reasoning-corruption & ML/AI & Xia 2026 \\
plan anchor & tool-output & reasoning-corruption & ML/AI & Yu 2026 \\
observation over-ingestion & tool-output & reasoning-corruption & ML/AI & Gaur 2026 \\
AdapTools & tool-output & goal-hijack & Security & Wang 2026 \\
Adversarial HalluSquatting & tool-output & goal-hijack & Security & Spira 2026 \\
Agent Data Injection (ADI) & tool-output & goal-hijack & Security & Choi 2026 \\
Agent-as-a-Proxy attack & tool-output & goal-hijack & Security & Isbarov 2026 \\
AGENTBAIT & tool-output & goal-hijack & Security & Wu 2026 \\
AgentVigil & tool-output & goal-hijack & Security & Wang 2025 \\
AIShellJack (coding-rule/config-file prompt injection against agentic IDE editors) & tool-output & goal-hijack & Security & Liu 2025 \\
ASPI (Ambiguous-State Prompt Injection) & tool-output & goal-hijack & Security & Sehwag 2026 \\
AutoDojo (adaptive black-box IPI attack generation) & tool-output & goal-hijack & Security & Ma 2026 \\
AutoInject & tool-output & goal-hijack & Security & Chen 2026 \\
autonomy tax & tool-output & reasoning-corruption & Security & Li 2026 \\
AutoRed & tool-output & goal-hijack & Security & Wang 2024 \\
BIPIA benchmark attack set & tool-output & goal-hijack & Security & Yi 2023 \\
ChainFuzzer & tool-output & goal-hijack & Security & Wu 2026 \\
Checkpoint-GCG & tool-output & goal-hijack & Security & Yang 2025 \\
ClawTrojan multi-step trojan/backdoor attack paradigm & tool-output & persistence-backdoor & Security & Tan 2026 \\
COMA & tool-output & goal-hijack & Security & Liu 2025 \\
Cross-App Context Poisoning (Confused ChatGPT) & tool-output & goal-hijack & Security & Wang 2026 \\
DataFlip & tool-output & goal-hijack & Security & Choudhary 2025 \\
EchoLeak (zero-click IPI, CVE-2025-32711) & tool-output & data-exfiltration & Security & Reddy 2025 \\
F2A (Feign Agent Attack) & tool-output & goal-hijack & Security & Ren 2024 \\
FlashRT & tool-output & goal-hijack & Security & Wang 2026 \\
Foot-in-the-Door (FITD) attack & tool-output & goal-hijack & Security & Nakash 2024 \\
Fun-tuning & tool-output & goal-hijack & Security & Labunets 2025 \\
HARE (Hybrid Agent Roleplay Exploitation) / Token Economy Poisoning & tool-output & goal-hijack & Security & Muhil 2025 \\
Indirect Prompt Injection (IPI) & tool-output & goal-hijack & Security & Greshake 2023 \\
Inter-Agent Trust Exploitation (AI agent privilege escalation) & tool-output & goal-hijack & Security & Lupinacci 2025 \\
IterInject & tool-output & goal-hijack & Security & Chen 2026 \\
Log-Substrate Prompt Injection & tool-output & goal-hijack & Security & Pandey 2026 \\
LogInject / Passive Prompt Injection & tool-output & goal-hijack & Security & Karanjai 2026 \\
LogJack & tool-output & goal-hijack & Security & Shah 2026 \\
MUZZLE & tool-output & goal-hijack & Security & Syros 2026 \\
Neural Exec & tool-output & goal-hijack & Security & Pasquini 2024 \\
ObliInjection & tool-output & goal-hijack & Security & Wang 2025 \\
Parallel-Poisoned Web (fingerprinting + website cloaking against AI web agents) & tool-output & goal-hijack & Security & Zychlinski 2025 \\
PI-Hunter & tool-output & goal-hijack & Security & He 2026 \\
PINA & tool-output & goal-hijack & Security & Liu 2026 \\
PISmith & tool-output & goal-hijack & Security & Yin 2026 \\
Prompt Overflow Attack & tool-output & goal-hijack & Security & Zhou 2026 \\
PromptFuzz & tool-output & goal-hijack & Security & Yu 2024 \\
PromptFuzz-SC & tool-output & goal-hijack & Security & Ren 2026 \\
query-based black-box optimization via LLM-expressed confidence (binary-comparison hill climbing) & tool-output & goal-hijack & Security & Zhang 2025 \\
Reasoning-Extension Denial-of-Service (DoS) & tool-output & resource-abuse & Security & Zhou 2026 \\
Relinking & tool-output & goal-hijack & Security & Liu 2026 \\
RL-Hammer & tool-output & goal-hijack & Security & Wen 2025 \\
RTA-PostForge (response-path tampering attack) & tool-output & goal-hijack & Security & Luo 2026 \\
SafeSearch & tool-output & goal-hijack & Security & Dong 2025 \\
Semantic-Masquerading IPI; Cascading IPI; Isolation-Breach IPI & tool-output & goal-hijack & Security & Ji 2025 \\
Silent Egress (implicit prompt injection via URL preview/metadata) & tool-output & data-exfiltration & Security & Lan 2026 \\
StruPhantom & tool-output & goal-hijack & Security & Feng 2025 \\
Targeted Promptware Attacks & tool-output & goal-hijack & Security & Nassi 2025 \\
ToolLeak & tool-output & goal-hijack & Security & Xie 2025 \\
TopicAttack & tool-output & goal-hijack & Security & Chen 2025 \\
\emph{ChatGPT memory/URL-exfiltration prompt injection bypass} & tool-output & data-exfiltration & Security & Schwartzman 2024 \\
\emph{data flow-based injection attack for observed-data exfiltration ('Important message' template)} & tool-output & data-exfiltration & Security & Alizadeh 2025 \\
\emph{HTML accessibility-tree universal adversarial trigger injection (TWUI)} & tool-output & goal-hijack & Security & Johnson 2025 \\
\emph{in-browser LLM-guided fuzzing for agentic-browser IPI} & tool-output & goal-hijack & Security & Cohen 2025 \\
\emph{log-interpretation evasion prompt injection (malicious-as-benign misclassification)} & tool-output & goal-hijack & Security & Landauer 2026 \\
\emph{multi-probe evasive GCG suffix attack} & tool-output & goal-hijack & Security & Rahman 2026 \\
\emph{PI-based counter-attack on offensive AI cybersecurity agents (four-stage attack anatomy + seven-category taxonomy incl. Unicode homograph exploitation)} & tool-output & goal-hijack & Security & Vilches 2025 \\
\emph{privacy leakage attack chain (external-content deception + IPI + exfiltration)} & tool-output & data-exfiltration & Security & Yang 2026 \\
VeriGrey (grey-box agent fuzzing for IPI) & tool-output & goal-hijack & Security & Zhang 2026 \\
VortexPIA & tool-output & data-exfiltration & Security & Cui 2025 \\
ZombAI & tool-output & persistence-backdoor & Security & Bayu 2026 \\
distraction by irrelevant context (graded by semantic relatedness) & RAG & reasoning-corruption & ML/AI & Wu 2024 \\
identity drift (persona coherence decay in generative agents) & RAG & reasoning-corruption & ML/AI & Platnick 2025 \\
context compliance (failure to track whether output was actually controlled by retrieved context vs. parametric prior vs. post-hoc rationalization) & RAG & silent-corruption & ML/AI & Chen 2026 \\
mechanistic explanation of knowledge conflict via information-flow/attention-head circuit (PH3) & RAG & silent-corruption & ML/AI & Jin 2024 \\
\emph{over-trust bias toward conflicting external evidence over correct parametric memory in RALMs} & RAG & silent-corruption & ML/AI & Jin 2024 \\
knowledge conflict (contextual vs parametric) & RAG & reasoning-corruption & ML/AI & Yuan 2024 \\
AGGD (Approximate Greedy Gradient Descent) & RAG & silent-corruption & Security & Su 2024 \\
BadRAG & RAG & goal-hijack & Security & Xue 2024 \\
ConfusedPilot & RAG & data-exfiltration & Security & RoyChowdhury 2024 \\
DACSI (Document-Authored Control-Signal Impersonation) & RAG & goal-hijack & Security & Zhu 2026 \\
Ignore; Detect; Exploit (hidden prompt injection attacks) & RAG & goal-hijack & Security & Collu 2025 \\
IKEA (Implicit Knowledge Extraction Attack) & RAG & data-exfiltration & Security & Wang 2025 \\
Indirect Data Poisoning (Distributed Denial of Science) & RAG & silent-corruption & Security & Gyevnár 2026 \\
Injection Paradox & RAG & goal-hijack & Security & Paeng 2026 \\
LogicPoison & RAG & goal-hijack & Security & Xiao 2026 \\
NeuroGenPoisoning & RAG & goal-hijack & Security & Zhu 2025 \\
Oracle Poisoning & RAG & reasoning-corruption & Security & Kereopa-Yorke 2026 \\
PANDORA (RAG Poisoning jailbreak) & RAG & goal-hijack & Security & Deng 2024 \\
PIDP-Attack & RAG & goal-hijack & Security & Wang 2026 \\
PoisonedRAG & RAG & goal-hijack & Security & Zou 2024 \\
Prompt-Injected Data Extraction & RAG & data-exfiltration & Security & Qi 2024 \\
RefineRAG & RAG & goal-hijack & Security & Wang 2026 \\
Salience Induction & RAG & goal-hijack & Security & Zhou 2026 \\
TrojanRAG & RAG & persistence-backdoor & Security & Cheng 2024 \\
\emph{backdoor-based RAG data extraction attack} & RAG & data-exfiltration & Security & Peng 2024 \\
\emph{gradient-based corpus poisoning (adversarial passage injection)} & RAG & silent-corruption & Security & Zhong 2023 \\
\emph{RAG-manipulation attack on multi-agent debate consensus} & RAG & reasoning-corruption & Security & Kraidia 2026 \\
\emph{retrieval-aware IPI (trigger-fragment decomposition attack)} & RAG & goal-hijack & Security & Chang 2026 \\
\emph{retriever fine-tuning backdoor attack} & RAG & persistence-backdoor & Security & Clop 2024 \\
\emph{tree-of-attacks ranking manipulation for conversational search} & RAG & goal-hijack & Security & Pfrommer 2024 \\
unintended long-term state poisoning & memory & silent-corruption & Both & Xu 2026 \\
multi-target memory interference & memory & silent-corruption & ML/AI & Lee 2026 \\
\emph{temporally-distant information inaccessibility in long-horizon agentic reasoning} & memory & reasoning-corruption & ML/AI & Hu 2026 \\
memory pollution (agent long-term memory) & memory & silent-corruption & ML/AI & Liu 2026 \\
memory laundering & memory & silent-corruption & ML/AI & Wang 2026 \\
memory misevolution & memory & silent-corruption & ML/AI & Xie 2026 \\
\emph{"cognitive event horizon" (context-saturation threshold beyond which accuracy collapses)} & memory & reasoning-corruption & ML/AI & Milošević 2026 \\
\emph{granularity mismatch in agent memory retrieval (instance-level memory causing both irrelevant-guidance retrieval and missed cross-task reuse)} & memory & reasoning-corruption & ML/AI & Shen 2026 \\
AgentPoison & memory & persistence-backdoor & Security & Chen 2024 \\
Chain-of-Trigger Backdoor (CoTri) & memory & persistence-backdoor & Security & Qiu 2025 \\
Cliopatra & memory & data-exfiltration & Security & Annamalai 2026 \\
Cross-Session Stored Prompt Injection (XSPI) & memory & persistence-backdoor & Security & Xie 2026 \\
MemoryGraft & memory & persistence-backdoor & Security & Srivastava 2025 \\
MemVenom & memory & persistence-backdoor & Security & Zhang 2026 \\
Plan Injection & memory & goal-hijack & Security & Patlan 2025 \\
Sleeper Memory Poisoning (Actor-Critic universal poisoning attack) & memory & persistence-backdoor & Security & Pulipaka 2026 \\
temporal memory contamination & memory & persistence-backdoor & Security & Al-Tawaha 2026 \\
\emph{memory-poisoning attack taxonomy (6 classes incl. Conditional Command Insertion, Policy-Conformant Fact Injection, False Precedent Insertion)} & memory & persistence-backdoor & Security & Dash 2026 \\
Zombie Agent & memory & persistence-backdoor & Security & Yang 2026 \\
context unawareness (failure to identify a small decisive piece of evidence within long/complex context) & cross-modal & reasoning-corruption & ML/AI & Xu 2026 \\
evidence misalignment (agentic long-video QA) & cross-modal & silent-corruption & ML/AI & Qiu 2026 \\
cross-modality parametric knowledge conflict & cross-modal & silent-corruption & ML/AI & Zhu 2024 \\
multimodal retrieval heads (sparse heads carrying retrieval) & cross-modal & reasoning-corruption & ML/AI & Li 2026 \\
AgentTypo & cross-modal & goal-hijack & Security & Li 2025 \\
AudioHijack & cross-modal & goal-hijack & Security & Chen 2026 \\
CrossInject & cross-modal & goal-hijack & Security & Wang 2025 \\
Graph-based prompt injection (math-expression-to-graph encoding) & cross-modal & goal-hijack & Security & Kwon 2024 \\
Image-based Prompt Injection (IPI) & cross-modal & goal-hijack & Security & Nagaraja 2025 \\
incidental prompt injection & cross-modal & reasoning-corruption & Security & Clusmann 2024 \\
Mind Mapping Prompt Injection & cross-modal & goal-hijack & Security & Lee 2025 \\
MIRAGE & cross-modal & goal-hijack & Security & Guo 2026 \\
MIRAGE (Masked Injection via Residual And Generative Exploitation) & cross-modal & goal-hijack & Security & Dai 2026 \\
Physical Prompt Injection Attack (PPIA) & cross-modal & goal-hijack & Security & Ling 2026 \\
PI3D & cross-modal & goal-hijack & Security & Li 2026 \\
self-interpreting images via adversarial meta-instructions & cross-modal & silent-corruption & Security & Zhang 2024 \\
steganographic prompt injection (multi-domain spatial/frequency/neural embedding) & cross-modal & goal-hijack & Security & Pathade 2025 \\
\emph{multi-modal prompt injection against robot navigation} & cross-modal & goal-hijack & Security & Zhang 2024 \\
WebInject & cross-modal & goal-hijack & Security & Wang 2025 \\
context rot & tool-metadata & reasoning-corruption & ML/AI & Ayyagari 2026 \\
Attractive Metadata Attack (AMA) & tool-metadata & goal-hijack & Security & Mo 2025 \\
AutoMalTool & tool-metadata & goal-hijack & Security & He 2025 \\
Branded Whisper Attack / Vault Whisper Attack (from 'Whispers of Wealth') & tool-metadata & resource-abuse & Security & Debi 2026 \\
ChatInject & tool-metadata & goal-hijack & Security & Chang 2025 \\
Cross-Tool Description Poisoning & tool-metadata & goal-hijack & Security & Shi 2026 \\
MCP-ITP (Implicit Tool Poisoning) & tool-metadata & goal-hijack & Security & Li 2026 \\
Sampling-Based Injection (MCP sampling/createMessage attack) & tool-metadata & goal-hijack & Security & Maloyan 2026 \\
TIP (Tree-structured Injection for Payloads) / stealthy update attack & tool-metadata & goal-hijack & Security & Shen 2026 \\
ToolHijacker & tool-metadata & goal-hijack & Security & Shi 2025 \\
\emph{descriptor-level semantic manipulation attack (MCP tool-selection bias)} & tool-metadata & goal-hijack & Security & Jamshidi 2025 \\
VATS (Error-Path Injection via Implicit Authority) & tool-metadata & goal-hijack & Security & Patel 2026 \\
extraneous cognitive load (context engineering overload) & multi-agent & reasoning-corruption & ML/AI & Shang 2025 \\
Constraint Drift (under long tool traces) & multi-agent & reasoning-corruption & ML/AI & Bui 2026 \\
context poisoning (self-induced, non-adversarial) & multi-agent & reasoning-corruption & ML/AI & Chen 2026 \\
MAD-Spear (Conformity-Driven Prompt Injection Attack) & multi-agent & goal-hijack & Security & Cui 2025 \\
MAS hijacking & multi-agent & goal-hijack & Security & Triedman 2025 \\
OMNI-LEAK (Orchestrator Multi-Agent Network Induced Data Leakage) & multi-agent & data-exfiltration & Security & Naik 2026 \\
PlanFlip & multi-agent & goal-hijack & Security & Wang 2026 \\
Prompt Infection & multi-agent & goal-hijack & Security & Lee 2025 \\
\emph{subagent inheritance exploitation (4 classes: insecure memory inheritance, absent resource access control, async memory divergence, unauthorized cross-agent termination)} & multi-agent & goal-hijack & Security & Cai 2026 \\
backdoor-powered prompt injection attack & supply-chain & persistence-backdoor & Security & Chen 2025 \\
Parasitic Toolchain Attacks (MCP Unintended Privacy Disclosure, MCP-UPD) & supply-chain & goal-hijack & Security & Zhao 2025 \\
PoisonedAlign & supply-chain & goal-hijack & Security & Shao 2024 \\
PoisonedSkills / DDIPE (Document-Driven Implicit Payload Execution) & supply-chain & goal-hijack & Security & Qu 2026 \\
\emph{chat-template inference-time backdoor} & supply-chain & persistence-backdoor & Security & Fogel 2026 \\
Virtual Prompt Injection (VPI) & supply-chain & persistence-backdoor & Security & Yan 2023 \\
\emph{agent-skill requirement decay under long trajectories (skill/requirement fade)} & skill & reasoning-corruption & ML/AI & Xue 2026 \\
POISE & skill & goal-hijack & Security & Hao 2026 \\
SkillJect & skill & goal-hijack & Security & Jia 2026 \\
\emph{skill-file prompt injection (trivial/non-optimized variant)} & skill & goal-hijack & Security & Schmotz 2025 \\
\end{longtable}
\endgroup


% GENERATED by scripts/gen_census_appendix.py -- do not hand-edit.
\begingroup
\footnotesize
\setlength{\tabcolsep}{3.5pt}
\begin{longtable}{@{}p{0.30\textwidth}lllll@{}}
\caption{Every named defense in the corpus (541 entries). \emph{Validated vs} is the threat model the paper's own evaluation tested, not the one it claims. Conventions otherwise as in Table~\ref{tab:census-mech}.}\\
\toprule
\textbf{Defense or fix} & \textbf{Channel} & \textbf{Acts at} & \textbf{Track} & \textbf{Validated vs} & \textbf{Source} \\
\midrule
\endfirsthead
\multicolumn{6}{l}{\emph{Table~\ref{tab:census-def}: named defenses, continued}} \\
\toprule
\textbf{Defense or fix} & \textbf{Channel} & \textbf{Acts at} & \textbf{Track} & \textbf{Validated vs} & \textbf{Source} \\
\midrule
\endhead
\midrule
\multicolumn{6}{r}{\emph{continued}} \\
\endfoot
\bottomrule
\endlastfoot
ActQKV (Activation-aware Probe-Query) & direct-input & reasoning & safety & incidental & \cite{xiao-2025-activationaware} \\
Agent-BRACE (Agent Belief state Representation via Abstraction and Confidence Estimation) & direct-input & none & safety & incidental & \cite{singh-2026-agentbrace} \\
Chain-of-Memory (CoM) - verifiable retrieval + causal tracing & direct-input & reasoning & safety & incidental & \cite{jarsania-2025-chainofmemory} \\
CoCoA (Confidence- and Context-Aware Adaptive Decoding) & direct-input & reasoning & safety & incidental & \cite{khandelwal-2025-cocoa} \\
CoThinker & direct-input & reasoning & safety & incidental & \cite{shang-2025-united} \\
Differential Transformer (DIFF Transformer) & direct-input & reasoning & safety & incidental & \cite{ye-2024-differential} \\
Ebbinghaus Forgetting Curve Memory Enhancement Framework & direct-input & reasoning & safety &  & \cite{jiang-2026-ebbinghaus} \\
EvoDS - Adaptive Context Compression (ACC) component specifically & direct-input & none & safety & incidental & \cite{yang-2026-evods} \\
FinLongDocAgent & direct-input & none & safety & incidental & \cite{wang-2026-documentlevel} \\
Free()LM (Free-Module LoRA adapter) & direct-input & reasoning & safety & incidental & \cite{zheng-2026-free} \\
G-MemLLM & direct-input & reasoning & safety & incidental & \cite{xu-2026-gmemllm} \\
GCAD (Gated Cropped Attention-Delta steering) & direct-input & reasoning & safety & incidental & \cite{kang-2026-promptactivation} \\
GDWM (Gated Differentiable Working Memory) & direct-input & reasoning & safety & incidental & \cite{mei-2026-gated} \\
JUICE (Just Run Twice) & direct-input & reasoning & safety & incidental & \cite{li-2025-taming} \\
KScope-guided Constrained Context Summarization & direct-input & none & safety & incidental & \cite{xiao-2025-kscope} \\
L2A (Learning To Attend) & direct-input & reasoning & safety & incidental & \cite{choudhary-2026-learning} \\
lambda-RLM & direct-input & reasoning & safety & incidental & \cite{roy-2026-ycombinator} \\
LCLM (Latent Context Language Models) & direct-input & ingestion & safety & incidental & \cite{li-2026-endtoend} \\
LLM+DPP (Principled Content Selection pipeline) & direct-input & ingestion & safety & incidental & \cite{padmakumar-2025-principled} \\
LongCrafter & direct-input & reasoning & safety & incidental & \cite{yuan-2026-longcrafter} \\
LOOC (Leave-One-Out Compression) & direct-input & ingestion & safety & incidental & \cite{wang-2026-alleviating} \\
P-RoPE / MiniWin (Periodic RoPE) & direct-input & reasoning & safety & incidental & \cite{huo-2026-periodic} \\
Planning with Schemas (PwS) & direct-input & reasoning & safety & incidental & \cite{xu-2025-cognitive} \\
PRIMA (context-boundary discipline + reference-material sandboxing + poisoned-output rejection) & direct-input & execution & safety & adversarial & \cite{annapureddy-2026-prima} \\
Program-Guided (PG) Runtime / Compile-then-Page SOP compiler & direct-input & execution & safety & incidental & \cite{yu-2026-compile} \\
QUITO-X & direct-input & ingestion & safety & incidental & \cite{wang-2024-quitox} \\
QwenLong-CPRS & direct-input & ingestion & safety & incidental & \cite{shen-2025-qwenlongcprs} \\
Recursive Models (recursive self-invocation with isolated subtask contexts) & direct-input & reasoning & safety & incidental & \cite{yang-2026-recursive} \\
Retrieve-then-Solve & direct-input & none & safety & incidental & \cite{du-2025-context} \\
RL Belief-State Reward Training + Representation-Level Steering (CBM mitigation) & direct-input & reasoning & safety & incidental & \cite{xu-2026-should} \\
SAFARI & direct-input & none & safety &  & \cite{zhu-2026-safari} \\
Scale Positional Hidden States (single-dimension scaling) & direct-input & reasoning & safety & incidental & \cite{yu-2024-mitigate} \\
State-Update Multi-turn Dialogue Strategy & direct-input & reasoning & safety & incidental & \cite{liu-2025-stateupdate} \\
Still (amortized KV cache compaction, Perceiver-based) & direct-input & ingestion & safety & incidental & \cite{oneill-2026-still} \\
TAG (Tagging-Augmented Generation) & direct-input & ingestion & safety & incidental & \cite{pal-2025-taggingaugmented} \\
TRIM-KV & direct-input & ingestion & safety & incidental & \cite{bui-2025-cache} \\
\emph{two preemptive-answer mitigations} & direct-input & reasoning & safety & both & \cite{xu-2024-preemptive} \\
AIR (Augmented Intermediate Representations) & direct-input & reasoning & security & adversarial & \cite{kariyappa-2025-stronger} \\
ARSM-Agent & direct-input & execution & security & adversarial & \cite{hu-2026-research} \\
ASIDE (Architecturally Separated Instruction-Data Embeddings) & direct-input & reasoning & security & adversarial & \cite{zverev-2025-aside} \\
Attack Detector Algorithm and Attack Library & direct-input & ingestion & security &  & \cite{alla-2025-cyberattacks} \\
BinaryShield & direct-input & ingestion & security & adversarial & \cite{gill-2025-binaryshield} \\
CaptureGuard & direct-input & ingestion & security & adversarial & \cite{kholkar-2025-capture} \\
ChatTEDU dual-LLM architecture (LLM-1 moderation gatekeeper + LLM-2 response generator) & direct-input & ingestion & security & adversarial & \cite{emekci-2025-securing} \\
Chunk-based hash + embedding prompt injection detection scheme (with optional blockchain audit logging) & direct-input & ingestion & security & adversarial & \cite{hao-2025-privacypreserving} \\
CoMAI & direct-input & none & security &  & \cite{sun-2026-comai} \\
Contiguity-gate defense (for guardrail sliding-window aggregation) & direct-input & ingestion & security & adversarial & \cite{zhou-2026-prompt} \\
data defenses & direct-input & ingestion & security & adversarial & \cite{agnew-2024-data} \\
DataSentinel & direct-input & ingestion & security & adversarial & \cite{liu-2025-datasentinel} \\
Defe-Prompt & direct-input & none & security & adversarial & \cite{ren-2024-f2a} \\
DefensiveToken (test-time optimized special tokens) & direct-input & ingestion & security & adversarial & \cite{chen-2025-defending} \\
DMPI-PMHFE & direct-input & ingestion & security & adversarial & \cite{ji-2025-detection} \\
DP-Fusion & direct-input & reasoning & security & adversarial & \cite{thareja-2025-dpfusion} \\
DRIP & direct-input & reasoning & security & adversarial & \cite{liu-2025-drip} \\
Enterprise Secure AI Execution Engine (Zero-Trust AI Security Gateway) & direct-input & ingestion & security & adversarial & \cite{krishna-2026-enterprise} \\
GenTel-Shield & direct-input & ingestion & security & adversarial & \cite{li-2024-gentelsafe} \\
GNN-based prompt injection detector & direct-input & ingestion & security & adversarial & \cite{jadhav-2025-novel} \\
Governed Execution (modular prompt orchestration for RAG-grounded call-center agents) & direct-input & none & security & both & \cite{pachtrachai-2026-transcripts} \\
GPT-OSSafeGuard + structured-representations + GSPO defense pipeline & direct-input & reasoning & security & adversarial & \cite{gudinorosero-2025-prompt} \\
Hybrid Constitutional Classifier & direct-input & ingestion & security & adversarial & \cite{lan-2025-hybrid} \\
IDE-Sanitizer & direct-input & ingestion & security & adversarial & \cite{nkoro-2026-idesanitizer} \\
IH-Challenge (RL training dataset with online adversarial example generation) & direct-input & reasoning & security & adversarial & \cite{guo-2026-ihchallenge} \\
IHDec (Instruction Hierarchy-steered Decoding) & direct-input & reasoning & security & both & \cite{liu-2026-ihdec} \\
InstruCoT & direct-input & reasoning & security & adversarial & \cite{chang-2026-know} \\
Instruction Hierarchy (privileged-instruction training) & direct-input & reasoning & security & adversarial & \cite{wallace-2024-instruction} \\
ISE (Instructional Segment Embedding) & direct-input & reasoning & security & adversarial & \cite{wu-2024-instructional} \\
Jatmo & direct-input & reasoning & security & adversarial & \cite{piet-2023-jatmo} \\
LIMBOSH & direct-input & execution & security &  & \cite{johnson-2024-modular} \\
LLM invocation logging & direct-input & execution & security &  & \cite{chernyshev-2023-large} \\
LLM-as-a-Judge (taxonomy-guided structured-reasoning judge) with Mixture-of-Models aggregation & direct-input & ingestion & security & adversarial & \cite{le-2026-prompt} \\
LLMON (LLM Object Notation) & direct-input & ingestion & security & adversarial & \cite{hind-2026-llmon} \\
LPDS (LLM Prompt Detection System) & direct-input & ingestion & security & adversarial & \cite{kim-2024-protection} \\
MAGE (Memory As Guardrail Enforcement) & direct-input & execution & security & adversarial & \cite{wang-2026-mage} \\
Meticulous Thought Defender & direct-input & reasoning & security & adversarial & \cite{shi-2025-meticulous} \\
Polymorphic Prompt Assembling (PPA) & direct-input & ingestion & security & adversarial & \cite{wang-2025-protect} \\
Prompt Injection as a Defense (against LLM-based PIE) & direct-input & ingestion & security & adversarial & \cite{liu-2024-evaluating} \\
Prompt Injection Threat Analysis + LLM Output Risk Monitoring framework & direct-input & ingestion & security & adversarial & \cite{zhao-2025-novel} \\
Prompt-response verification framework (Task Analyst + Safety Analyst + Judge) & direct-input & reasoning & security & adversarial & \cite{thota-2026-verifying} \\
Prompt-Shield Framework (PSF) & direct-input & ingestion & security & adversarial & \cite{hadiprakoso-2025-adaptive} \\
PromptGuard & direct-input & ingestion & security & adversarial & \cite{alzahrani-2026-promptguard} \\
PromptLocate (injected-prompt localization) & direct-input & ingestion & security & adversarial & \cite{jia-2025-promptlocate} \\
PromptSleuth & direct-input & ingestion & security & adversarial & \cite{wang-2025-promptsleuth} \\
Route-Safety Audit Contract (confirmation-plus-provenance routing) & direct-input & execution & security & adversarial & \cite{tai-2026-brainprompt} \\
SE+CFE (Standard Evaluation + Counterfactual Evaluation) & direct-input & reasoning & security & adversarial & \cite{liu-2025-counterfactual} \\
SecInfer & direct-input & reasoning & security & adversarial & \cite{liu-2025-secinfer} \\
Secure Artificial Intelligence (SAI) & direct-input & ingestion & security & adversarial & \cite{teckani-2025-secure} \\
Secure Prompt Engineering Patterns (SPEP) & direct-input & reasoning & security & adversarial & \cite{sudarsan-2026-secure} \\
SecureCAI & direct-input & reasoning & security & adversarial & \cite{ali-2026-securecai} \\
Sentinel (qualifire/prompt-injection-sentinel) & direct-input & ingestion & security & adversarial & \cite{ivry-2025-sentinel} \\
ShadowPlay & direct-input & ingestion & security & both & \cite{alsobeh-2025-shadowplay} \\
SHIELD & direct-input & ingestion & security & adversarial & \cite{esugo-2025-shield} \\
soft begging & direct-input & reasoning & security &  & \cite{ostermann-2024-soft} \\
SPIN (Self-Supervised Prompt INjection) & direct-input & reasoning & security & adversarial & \cite{zhou-2024-spin} \\
SPRA (Scaffold-Preserving Representation Alignment) & direct-input & reasoning & security & adversarial & \cite{shao-2026-mitigating} \\
StruQ (structured queries) & direct-input & ingestion & security & adversarial & \cite{chen-2024-struq} \\
TIVS-O agentic defense pipeline (Front-End Agent + Guard-Sanitizer + Policy Enforcer with semantic caching) & direct-input & reasoning & security & adversarial & \cite{gosmar-2026-prompt} \\
unified Prompt Assembling + State Management + Safety Validation framework & direct-input & execution & security & adversarial & \cite{zhang-2025-enhancing} \\
UniGuardian & direct-input & ingestion & security & adversarial & \cite{lin-2025-uniguardian} \\
\emph{domain-specific multi-layer safeguard pipeline (pattern filters + structural validation + contextual sandboxing + session-level behavioral checks)} & direct-input & reasoning & security & adversarial & \cite{maiorano-2026-evaluating} \\
\emph{Embedding-based traditional-ML classifiers for PI detection} & direct-input & ingestion & security & adversarial & \cite{ayub-2024-embeddingbased} \\
\emph{Explanation-based secondary detection signal for log-injection evasion} & direct-input & none & security & adversarial & \cite{landauer-2026-just} \\
\emph{Hybrid multi-layered real-time PI detection framework} & direct-input & ingestion & security & adversarial & \cite{prakash-2026-hybrid} \\
\emph{Hybrid SIEM + One-Class SVM prompt-injection detection framework} & direct-input & ingestion & security & adversarial & \cite{alshammari-2026-detecting} \\
\emph{ICL-based multi-LLM agent harmfulness-detection framework} & direct-input & reasoning & security & adversarial & \cite{sato-2026-addressing} \\
\emph{Multi-agent PI defense pipeline (Chain-of-Agents / Coordinator-based)} & direct-input & reasoning & security & adversarial & \cite{hossain-2025-multiagent} \\
\emph{multi-agent privacy firewall (browser extension + MITM proxy + PII-detection pipeline)} & direct-input & ingestion & security & incidental & \cite{cuesta-2026-multiagent} \\
\emph{P2SQL mitigations (DB permission hardening + query rewriting)} & direct-input & execution & security & adversarial & \cite{pedro-2023-prompt} \\
\emph{Positive-Security Prompt Firewall for healthcare PI/PHI protection} & direct-input & ingestion & security & adversarial & \cite{schwoebel-2026-injection} \\
\emph{secure prompt strategies for robotic navigation} & direct-input & ingestion & security & adversarial & \cite{zhang-2024-study} \\
\emph{step-level fine-tuned industrial tool-call guardrail} & direct-input & execution & security & adversarial & \cite{xue-2026-multiagent} \\
\emph{Trust Context Separation for privileged AI agent systems} & direct-input & execution & security & adversarial & \cite{singh-2026-securingb} \\
\emph{Uncertainty-based Random Forest attack detector (chimera/Xmera)} & direct-input & none & security & adversarial & \cite{fastowski-2025-injecting} \\
ACON (Agent Context Optimization) & tool-output & ingestion & safety & incidental & \cite{kang-2025-acon} \\
ActiveContext (ContextCurator + TaskExecutor) & tool-output & ingestion & safety & incidental & \cite{li-2026-escaping} \\
AdaCoM (Adaptive Context Management) & tool-output & reasoning & safety & incidental & \cite{yi-2026-learning} \\
AgentSwing & tool-output & reasoning & safety & incidental & \cite{feng-2026-agentswing} \\
ARC (Addressable Recall Compaction) & tool-output & ingestion & safety & incidental & \cite{dang-2026-addressable} \\
BACM / BACM-RL (Budget-Aware Context Management) & tool-output & ingestion & safety & incidental & \cite{wu-2026-contextbudget} \\
ClinQueryAgent (Retrieval Agent sub-agent delegation) & tool-output & none & safety & incidental & \cite{boyle-2026-clinqueryagent} \\
CodeDelegator (role separation via Delegator/Coder agents + EPSS) & tool-output & execution & safety & incidental & \cite{fei-2026-codedelegator} \\
ColorBrowserAgent (knowledge-aligned progressive summarization + human-in-the-loop knowledge adaptation) & tool-output & none & safety & incidental & \cite{wang-2026-colorbrowseragent} \\
CompactionRL & tool-output & ingestion & safety & incidental & \cite{li-2026-compactionrl} \\
COMPASS (Context-Organized Multi-Agent Planning and Strategy System) & tool-output & reasoning & safety & incidental & \cite{wan-2025-compass} \\
Context Viewer & tool-output & none & safety &  & \cite{sriraman-2026-context} \\
Context-Folding (FoldGRPO) & tool-output & ingestion & safety & incidental & \cite{sun-2025-scaling} \\
FocusAgent / DefenseFocusAgent & tool-output & ingestion & safety & both & \cite{kerboua-2025-focusagent} \\
icat-agent (decentralized event-based multi-agent scaffold) & tool-output & none & safety & incidental & \cite{chen-2026-unlocking} \\
MAT-based Assurance Framework (Message-Action Traces + governance layer: capability limits + action mediation) & tool-output & execution & safety &  & \cite{paduraru-2026-tracebased} \\
PAACE (Plan-Aware Automated Context Engineering) & tool-output & ingestion & safety & incidental & \cite{yuksel-2025-paace} \\
PRInTS & tool-output & none & safety &  & \cite{lee-2025-prints} \\
Prometheus (memory-centric context engine with working memory over a repository knowledge graph) & tool-output & none & safety & incidental & \cite{pan-2025-prometheus} \\
ReSum / ReSum-GRPO & tool-output & ingestion & safety & incidental & \cite{wu-2025-resum} \\
SDO (Signal-Driven Observation) & tool-output & ingestion & safety & incidental & \cite{gaur-2026-signaldriven} \\
Self-GC (Self-Governing Context) & tool-output & ingestion & safety & incidental & \cite{hao-2026-selfgc} \\
Slipstream (asynchronous, judge-validated trajectory-grounded compaction) & tool-output & ingestion & safety & incidental & \cite{chen-2026-slipstream} \\
SmoothAgent & tool-output & ingestion & safety &  & \cite{pan-2026-smoothagent} \\
STRACE (Structural Trajectory Analysis and Causal Extraction) & tool-output & none & safety & incidental & \cite{chang-2026-noisy} \\
Synapse (trajectory-as-exemplar prompting + exemplar memory) & tool-output & none & safety & incidental & \cite{zheng-2023-synapse} \\
\emph{Rot-aware rejection-sampling filter for long-horizon search} & tool-output & reasoning & safety & incidental & \cite{xia-2026-diagnosing} \\
VISTA (Visible Internal State for Tool Agents) & tool-output & reasoning & safety & incidental & \cite{xu-2026-llm} \\
ACE (Abstract-Concrete-Execute) & tool-output & execution & security & adversarial & \cite{li-2025-ace} \\
Activation-based (linear probe) prompt-injection classifier, evaluated under Leave-One-Dataset-Out (LODO) protocol & tool-output & ingestion & security & adversarial & \cite{fomin-2026-benchmarks} \\
Agent Privilege Separation (Reader/Actor two-agent isolation + JSON formatting) & tool-output & execution & security & adversarial & \cite{cheng-2026-agent} \\
Agent Sentry (structural classifier + allowlist + LLM judge) & tool-output & execution & security & adversarial & \cite{sequeira-2026-agentsentry} \\
AgentArmor & tool-output & execution & security & adversarial & \cite{wang-2025-agentarmor} \\
Agentic Immune System (AIS) -- CIFT + DP-HONEY + cumulative leakage tracking & tool-output & reasoning & security & adversarial & \cite{chauhan-2026-caught} \\
AgenticOS (Ghost Kernel / Logic Shutter / Agent Capsule / Semantic Boundary Gateway) & tool-output & execution & security & adversarial & \cite{zhao-2026-agenticos} \\
AgentRedGuard & tool-output & none & security & adversarial & \cite{dingeto-2026-agentredbench} \\
AgenTRIM & tool-output & execution & security & adversarial & \cite{betser-2026-agentrim} \\
AgentSandbox & tool-output & none & security & adversarial & \cite{zhang-2025-llm} \\
AgentSentry & tool-output & ingestion & security & adversarial & \cite{zhang-2026-agentsentry} \\
AgentShield & tool-output & execution & security & adversarial & \cite{rassul-2026-agentshield} \\
AgentSight (eBPF boundary tracing) & tool-output & execution & security & both & \cite{zheng-2025-agentsight} \\
AgentSys & tool-output & ingestion & security & adversarial & \cite{wen-2026-agentsys} \\
AgentVisor & tool-output & execution & security & adversarial & \cite{ying-2026-agentvisor} \\
AgentWatcher & tool-output & ingestion & security & adversarial & \cite{wang-2026-agentwatcher} \\
aiAuthZ (identity-bound off-host authorization gateway) & tool-output & execution & security & adversarial & \cite{kodathala-2026-aiauthz} \\
AlignSentinel (three-class instruction-hierarchy-aware attention-based classifier) & tool-output & ingestion & security & adversarial & \cite{jia-2026-alignsentinel} \\
APPA (Agentic Permissions Policy Algebra) & tool-output & execution & security & adversarial & \cite{kravchenko-2026-agentic} \\
AprielGuard & tool-output & ingestion & security & adversarial & \cite{kasundra-2025-aprielguard} \\
ARGUS & tool-output & reasoning & security & adversarial & \cite{weng-2026-argus} \\
ARLAS (Adversarial Reinforcement Learning for Agent Safety) & tool-output & reasoning & security & adversarial & \cite{wang-2025-adversarial} \\
Attention Tracker & tool-output & reasoning & security & adversarial & \cite{hung-2024-attention} \\
AttriGuard & tool-output & execution & security & adversarial & \cite{he-2026-attriguard} \\
Authenticated Prompts / Authenticated Context (cryptographic provenance + policy algebra) & tool-output & execution & security & adversarial & \cite{rajagopalan-2026-protecting} \\
AuthGraph & tool-output & execution & security & adversarial & \cite{wang-2026-aligning} \\
Boundary awareness + explicit reminder & tool-output & ingestion & security & adversarial & \cite{yi-2023-benchmarking} \\
CaMeL & tool-output & execution & security & adversarial & \cite{debenedetti-2025-defeating} \\
ceLLMate & tool-output & execution & security & adversarial & \cite{meng-2025-cellmate} \\
ClawAudit & tool-output & none & security &  & \cite{zhang-2026-local} \\
ClawGuard & tool-output & execution & security & adversarial & \cite{zhao-2026-clawguard} \\
Cloud Decoy AI Agent & tool-output & execution & security &  & \cite{manivannan-2026-agentic} \\
CodeSentinel & tool-output & ingestion & security & adversarial & \cite{cheng-2026-codesentinel} \\
Cognitive Control Architecture (CCA) & tool-output & execution & security & adversarial & \cite{liang-2025-cognitive} \\
Cognitive Firewall & tool-output & ingestion & security & adversarial & \cite{lan-2026-cognitive} \\
CommandSans & tool-output & ingestion & security & adversarial & \cite{das-2025-commandsans} \\
CXI (Context-to-Execution Integrity) & tool-output & execution & security & adversarial & \cite{santosgrueiro-2026-contexttoexecution} \\
DataFilter & tool-output & ingestion & security & adversarial & \cite{wang-2025-defending} \\
DCP (Device Context Protocol) & tool-output & ingestion & security & adversarial & \cite{yang-2026-device} \\
DRIFT (Dynamic Rule-based Isolation Framework for Trustworthy agentic systems) & tool-output & execution & security & adversarial & \cite{li-2025-drift} \\
DualView & tool-output & ingestion & security & adversarial & \cite{kim-2026-dualview} \\
Encrypted Prompt & tool-output & execution & security &  & \cite{chan-2025-encrypted} \\
ESLD (External Surrogate Latent Defense) & tool-output & reasoning & security & adversarial & \cite{narendra-2026-esld} \\
f-secure LLM system (information flow control) & tool-output & execution & security & adversarial & \cite{wu-2024-systemlevel} \\
FASA (zero-trust, full-lifecycle defense philosophy) & tool-output & execution & security & adversarial & \cite{ying-2026-uncovering} \\
FATE (on-policy self-evolving repair framework) with Pareto-Front Policy Optimization (PFPO) & tool-output & reasoning & security & adversarial & \cite{yin-2026-onpolicy} \\
FATH (Formatting AuThentication with Hash-based tags) & tool-output & execution & security & adversarial & \cite{wang-2024-fath} \\
four layer defense in depth architecture & tool-output & execution & security &  & \cite{maiti-2026-caging} \\
GIF (Geometric Information Flow) & tool-output & reasoning & security & adversarial & \cite{storek-2026-gif} \\
GoalAnchorCheck (D1) / CrossAgentConsensus (D2) & tool-output & none & security & adversarial & \cite{wang-2026-planflip} \\
Hybrid context management + deterministic programmatic safety checks / domain allowlisting / action restriction / least-privilege scoping (production browser-agent architecture) & tool-output & execution & security & adversarial & \cite{vardanyan-2025-building} \\
ICON (Inference-Time Correction) & tool-output & reasoning & security & adversarial & \cite{wang-2026-icon} \\
InjectDefuser & tool-output & none & security & adversarial & \cite{koide-2026-clouding} \\
IntentGuard & tool-output & reasoning & security & adversarial & \cite{kang-2025-mitigating} \\
IPIGuard & tool-output & execution & security & adversarial & \cite{an-2025-ipiguard} \\
KBRA (Keyed Binding Reassembly Audit) & tool-output & ingestion & security & adversarial & \cite{liu-2026-safe} \\
LGA (Layered Governance Architecture) & tool-output & execution & security & adversarial & \cite{ge-2026-governance} \\
LlamaFirewall (PromptGuard 2 + AlignmentCheck + modular guardrail engine) & tool-output & ingestion & security & adversarial & \cite{chennabasappa-2025-llamafirewall} \\
LLMbda / Randori (dynamic information-flow-control lambda calculus with verified interpreter) & tool-output & reasoning & security & adversarial & \cite{garby-2026-llmbda} \\
Mantis & tool-output & none & security & adversarial & \cite{pasquini-2024-hacking} \\
MELON (Masked re-Execution and TooL comparisON) & tool-output & reasoning & security & adversarial & \cite{zhu-2025-melon} \\
MIPIAD & tool-output & reasoning & security & adversarial & \cite{muhtadi-2026-mipiad} \\
Mixture of Encodings & tool-output & ingestion & security & adversarial & \cite{zhang-2025-defense} \\
MOSAIC (plan-check-act/refuse post-training framework) & tool-output & reasoning & security & adversarial & \cite{agarwal-2026-learning} \\
NeuroTaint & tool-output & reasoning & security & adversarial & \cite{cai-2026-ghost} \\
OpenClaw PRISM & tool-output & execution & security & adversarial & \cite{li-2026-openclaw} \\
PACT (Provenance-Aware Capability Contracts) & tool-output & execution & security & adversarial & \cite{fan-2026-granularity} \\
Pattern D: Inference-boundary redaction proxy (PII extraction/tokenization before model inference) & tool-output & ingestion & security & adversarial & \cite{mohan-2026-agent} \\
PCFI (Prompt Control-Flow Integrity) & tool-output & reasoning & security & adversarial & \cite{alam-2026-prompt} \\
PIGuard (Mitigating Over-defense for Free) & tool-output & ingestion & security & adversarial & \cite{li-2025-piguard} \\
PIIGuard & tool-output & ingestion & security & adversarial & \cite{liu-2026-piiguard} \\
PISanitizer & tool-output & ingestion & security & adversarial & \cite{geng-2025-pisanitizer} \\
PIShield & tool-output & reasoning & security & adversarial & \cite{zou-2025-pishield} \\
Plan-Then-Execute (PTE) architecture for web agents & tool-output & execution & security & adversarial & \cite{piet-2026-web} \\
PlanGuard & tool-output & reasoning & security & adversarial & \cite{gong-2026-planguard} \\
PRISM (activation-conditioned instruction-set interpreter) & tool-output & reasoning & security & adversarial & \cite{gressel-2026-prism} \\
Progent & tool-output & execution & security & adversarial & \cite{shi-2025-progent} \\
ProjGuard & tool-output & execution & security & adversarial & \cite{contreras-2026-projguard} \\
prompt-shield (Seven Cross-Domain Detection Techniques) & tool-output & ingestion & security & adversarial & \cite{munirathinam-2026-pattern} \\
PromptArmor & tool-output & ingestion & security & adversarial & \cite{shi-2025-promptarmor} \\
PromptSentinel-X & tool-output & ingestion & security & adversarial & \cite{zhuhadar-2026-promptsentinelx} \\
PromptShield (fine-tuned deployable PI detector) & tool-output & ingestion & security & adversarial & \cite{jacob-2024-promptshield} \\
ReasAlign & tool-output & reasoning & security & adversarial & \cite{li-2026-reasalign} \\
RedVisor & tool-output & reasoning & security & adversarial & \cite{liu-2026-redvisor} \\
RETA (Reasoning-Enabled Task Alignment) & tool-output & reasoning & security & adversarial & \cite{he-2026-defending} \\
Robustness via Referencing (instruction-referencing defense) & tool-output & reasoning & security & adversarial & \cite{chen-2025-robustness} \\
RTBAS (Robust Tool-Based Agent System) & tool-output & execution & security & adversarial & \cite{zhong-2025-rtbas} \\
SafeAgent & tool-output & execution & security & adversarial & \cite{liu-2026-safeagent} \\
Safety-guided Chain-of-Thought (SCoT) & tool-output & none & security & both & \cite{lee-2024-mobilesafetybench} \\
SecAlign (preference optimization) & tool-output & reasoning & security & adversarial & \cite{chen-2024-secalign} \\
SecAlign++ / Meta-SecAlign & tool-output & reasoning & security & adversarial & \cite{chen-2025-meta} \\
SecMCP & tool-output & reasoning & security & adversarial & \cite{shi-2025-quantifying} \\
SecureClaw (dual-boundary: effect-sink authorization + read-boundary plaintext confinement) & tool-output & execution & security & adversarial & \cite{ma-2026-secureclaw} \\
SIC (Soft Instruction Control) & tool-output & reasoning & security & adversarial & \cite{walter-2025-soft} \\
Signed-Prompt & tool-output & ingestion & security & adversarial & \cite{suo-2024-signedprompt} \\
Spotlighting (delimiting / datamarking / encoding) & tool-output & ingestion & security & adversarial & \cite{hines-2024-defending} \\
SUPERVISOR & tool-output & none & security & adversarial & \cite{wu-2026-bots} \\
Task Shield & tool-output & reasoning & security & adversarial & \cite{jia-2024-task} \\
Tool Result Parsing (ParseData/CheckTool) & tool-output & ingestion & security & adversarial & \cite{yu-2026-defense} \\
ToolSafe (TS-Guard + TS-Flow) & tool-output & execution & security & adversarial & \cite{mou-2026-toolsafe} \\
Trace & tool-output & execution & security &  & \cite{ediga-2026-trace} \\
TRACE (Trajectory Risk-Aware Compression) & tool-output & ingestion & security & adversarial & \cite{hong-2026-trace} \\
TraceCaps & tool-output & execution & security & adversarial & \cite{catarino-2026-tracecaps} \\
TRACER-AI & tool-output & ingestion & security & adversarial & \cite{singh-2026-tracerai} \\
Twin Agent (Explore Agent / Safe Agent residual-hint privilege separation) & tool-output & execution & security & adversarial & \cite{hu-2026-twin} \\
Type-directed privilege separation & tool-output & execution & security & adversarial & \cite{jacob-2025-preventing} \\
\emph{Attack-derived shield-prompt defenses (Ignore/Escape/Fake-Completion/Combined Defense)} & tool-output & reasoning & security & adversarial & \cite{chen-2024-defense} \\
\emph{Authorization framework for autonomous LLM agent threat mitigation} & tool-output & none & security & adversarial & \cite{kota-2025-systemic} \\
\emph{Context-Aware Interception Framework for PI detection/mitigation} & tool-output & execution & security & adversarial & \cite{singh-2026-securing} \\
\emph{deterministic post-audit verifier layer added to a compositional safety system} & tool-output & none & security & adversarial & \cite{zhang-2026-ownerharm} \\
\emph{Digital forensic analysis framework for post-hoc IPI attack identification from agent execution logs} & tool-output & none & security & adversarial & \cite{chernyshev-2024-short} \\
\emph{Embedding-based semantic context analysis detector for IPI} & tool-output & ingestion & security & adversarial & \cite{alamsabi-2026-embeddingbased} \\
\emph{Four-layer guardrail system for offensive-security AI agents} & tool-output & none & security & adversarial & \cite{vilches-2025-cybersecurity} \\
\emph{Fraud-detection-inspired trajectory-feature detector (XGBoost, 42-dim features)} & tool-output & ingestion & security & adversarial & \cite{yu-2026-lowlatency} \\
\emph{IPI detection + post-detection removal pipeline} & tool-output & ingestion & security & adversarial & \cite{chen-2025-indirect} \\
\emph{Layered defense (input filtering + prompt hardening/spotlighting + output validation) against passive prompt injection} & tool-output & none & security & adversarial & \cite{karanjai-2026-context} \\
\emph{Policy Enforcement Point (PEP) for MCP tool-call boundary} & tool-output & execution & security & adversarial & \cite{wang-2026-runtime} \\
\emph{Pre-deployment scan/harden/validate pipeline (Analyzer/Modifier/Validator)} & tool-output & ingestion & security & adversarial & \cite{shukla-2026-data} \\
\emph{Reflection-based defenses against Foot-in-the-Door IPI (Self-Reflection / Hesitation Reflection / Safe Reflection)} & tool-output & none & security & adversarial & \cite{nakash-2024-breaking} \\
\emph{Representation Engineering-based Probing for IPI detection} & tool-output & none & security & adversarial & \cite{zhu-2026-your} \\
\emph{Scientist Persona + Provenance Audit mitigations against indirect data poisoning} & tool-output & none & security & adversarial & \cite{gyevnar-2026-distributed} \\
\emph{six composable agent design patterns (e.g. Action-Selector) for prompt-injection-resistant architecture} & tool-output & execution & security & adversarial & \cite{beurerkellner-2025-design} \\
\emph{WAB Wasm-sandboxed capability-based authority mechanism (allowlisted API with explicit grant/revocation)} & tool-output & execution & security & adversarial & \cite{liu-2026-wab} \\
User-Tool Aligned Safety RL Framework & tool-output & reasoning & security & adversarial & \cite{sha-2025-agent} \\
VIGIL & tool-output & execution & security & adversarial & \cite{lin-2026-vigil} \\
Windowed HTML content filtering (128/64-token chunking + PromptGuard/DeBERTaV3 classification) & tool-output & ingestion & security & adversarial & \cite{dinu-2025-disrupting} \\
Zero-Shot Embedding Drift Detection (ZEDD) & tool-output & ingestion & security & adversarial & \cite{sekar-2026-zeroshot} \\
A-RAG (hierarchical retrieval interfaces) & RAG & none & safety & incidental & \cite{du-2026-arag} \\
ACD (Adaptive Contrastive Decoding) & RAG & reasoning & safety & incidental & \cite{kim-2024-adaptive} \\
Astute RAG & RAG & reasoning & safety & incidental & \cite{wang-2024-astute} \\
AttnComp (Attention-guided Context Compression) & RAG & ingestion & safety & incidental & \cite{luo-2025-attncomp} \\
BordaRAG (Borda Voting-based document selection) & RAG & reasoning & safety & incidental & \cite{li-2025-bordarag} \\
CARE-RAG (Conflict-Aware and Reliable Evidence for RAG) & RAG & reasoning & safety & incidental & \cite{chen-2025-rethinking} \\
CD2 (Conflict-Disentangle Contrastive Decoding) & RAG & reasoning & safety & incidental & \cite{jin-2024-tugofwar} \\
CDD (Context-Driven Decomposition) & RAG & none & safety & adversarial & \cite{chen-2026-rag} \\
CheckRLM & RAG & reasoning & safety & incidental & \cite{xu-2026-checkrlm} \\
CK-PLUG (Controllable Knowledge Plug-in) & RAG & reasoning & safety & incidental & \cite{bi-2025-parameters} \\
CKG-RAG & RAG & reasoning & safety & incidental & \cite{zhang-2026-ckgrag} \\
COIECD (contextual information-entropy constraint decoding) & RAG & reasoning & safety & incidental & \cite{yuan-2024-discerning} \\
COMBO (Compatibility-Oriented knowledge Merging) & RAG & reasoning & safety & incidental & \cite{zhang-2023-merging} \\
Conflict-Aware RAG (ConScore signal + multi-stage SFT/DPO/reranking) & RAG & reasoning & safety & incidental & \cite{wu-2026-conflictaware} \\
ConflictRAG (two-stage detection + Entropy-TOPSIS credibility + CARS metric) & RAG & reasoning & safety & incidental & \cite{wang-2026-conflictrag} \\
Corrective Retrieval Augmented Generation (CRAG) & RAG & ingestion & safety & incidental & \cite{yan-2024-corrective} \\
Cross-Validation Re-ranking (inter-context conflict-aware re-ranking) & RAG & ingestion & safety & incidental & \cite{wu-2025-eliminating} \\
CRVA-TGRAG (Conflict Resolution in Vulnerability Analysis via Teacher-Guided RAG) & RAG & reasoning & safety & incidental & \cite{zhou-2026-tugofwar} \\
Discern-and-Answer & RAG & none & safety & incidental & \cite{hong-2023-gullible} \\
DualGraph & RAG & reasoning & safety & incidental & \cite{shi-2026-tale} \\
Dynamic Context Selection (classifier-based adaptive k-selection + reranking) & RAG & ingestion & safety & incidental & \cite{iratni-2025-dynamic} \\
DyPRAG (Dynamic Parametric RAG) & RAG & reasoning & safety & incidental & \cite{tan-2025-dynamic} \\
EHR-RAG & RAG & none & safety & incidental & \cite{cao-2026-ehrrag} \\
END (Early Noise Dropping) & RAG & ingestion & safety & incidental & \cite{jin-2025-end} \\
ERA (Evidence-based Reliability Alignment) & RAG & reasoning & safety & incidental & \cite{shin-2026-era} \\
FaithfulRAG (fact-level conflict modeling + Self-Think) & RAG & reasoning & safety & incidental & \cite{zhang-2025-faithfulrag} \\
GenDiE (Generate, Discriminate, Evolve) & RAG & reasoning & safety & incidental & \cite{li-2025-generate} \\
ITEM (Integrated Term Enhancement Methodology) & RAG & none & safety & incidental & \cite{bei-2024-manufacturing} \\
KCM (Knowledge Conflict Mitigation) & RAG & reasoning & safety & incidental & \cite{an-2026-enhancing} \\
Knowledgeable-R1 & RAG & reasoning & safety & incidental & \cite{lin-2025-resisting} \\
KnowPO (Knowledge-aware Preference Optimization) & RAG & reasoning & safety & incidental & \cite{zhang-2024-knowpo} \\
KV-Fusion & RAG & reasoning & safety & incidental & \cite{oh-2025-parallel} \\
KVEraser & RAG & reasoning & safety & incidental & \cite{li-2026-kveraser} \\
LDAR (Learning Distraction-Aware Retrieval) & RAG & ingestion & safety & incidental & \cite{shim-2025-rag} \\
LongFaith & RAG & reasoning & safety &  & \cite{yang-2025-longfaith} \\
LongRAG & RAG & reasoning & safety & incidental & \cite{zhao-2024-longrag} \\
MACR (multi-agent conflict resolution) & RAG & reasoning & safety & incidental & \cite{peng-2026-navigating} \\
MAL-RAG (Multiple Abstraction Level RAG) & RAG & ingestion & safety & incidental & \cite{zheng-2025-multiple} \\
Micro-Act (hierarchical action-space decomposition) & RAG & reasoning & safety & incidental & \cite{huo-2025-microact} \\
Mobile-Agent-RAG (Manager-RAG / Operator-RAG dual-level retrieval augmentation) & RAG & none & safety & incidental & \cite{zhou-2025-mobileagentrag} \\
ParamMute (Parametric Knowledge Muting through FFN Suppression) & RAG & reasoning & safety & incidental & \cite{huang-2025-parammute} \\
PH3 (Pruning Head via PatH PatcHing) & RAG & reasoning & safety & incidental & \cite{jin-2024-cutting} \\
ProbeRAG (knowledge pruning + latent conflict probing + conflict-aware attention) & RAG & reasoning & safety & incidental & \cite{gao-2025-blackbox} \\
RADIANT (DPO-ECD, Entity-Context Divergence alignment) & RAG & none & safety & incidental & \cite{rawte-2025-radiant} \\
RAPS-DA (Regime-Aware Peer Specialization) & RAG & reasoning & safety & both & \cite{wang-2026-regimeaware} \\
RAT (Retrieval Augmented Thoughts) & RAG & reasoning & safety & incidental & \cite{wang-2024-rat} \\
RE-IAG (Reasoning-Evidence Interactive Augmented Graph) & RAG & reasoning & safety & incidental & \cite{sui-2026-decomposition} \\
Reconstruct to Discriminate (RtD) & RAG & reasoning & safety & incidental & \cite{peng-2025-misinformation} \\
RPO (Retrieval Preference Optimization) & RAG & reasoning & safety & incidental & \cite{yan-2025-rpo} \\
SABER (Self-Aware Belief Estimator for RAG) & RAG & reasoning & safety & incidental & \cite{zhu-2026-trust} \\
SCR \& CR-DPO (Situated Faithfulness) & RAG & reasoning & safety & incidental & \cite{huang-2024-trust} \\
Self-RAG & RAG & reasoning & safety & incidental & \cite{asai-2023-selfrag} \\
SHIFT (gate-modulated activation steering) & RAG & reasoning & safety & incidental & \cite{li-2026-shift} \\
SpARE (Sparse Auto-Encoder Representation Engineering) & RAG & reasoning & safety & incidental & \cite{zhao-2024-steering} \\
SUGAR (Semantic Uncertainty Guided Adaptive Retrieval) & RAG & reasoning & safety & incidental & \cite{zubkova-2025-sugar} \\
Swin-VIB (variational information bottleneck pipeline) & RAG & reasoning & safety & incidental & \cite{wang-2025-accommodate} \\
Taxonomy-Aware / Pipeline conflict-type-aware response generation (from the CONFLICTS benchmark) & RAG & reasoning & safety & incidental & \cite{cattan-2025-dragged} \\
TCR (Transparent Conflict Resolution) & RAG & reasoning & safety & incidental & \cite{ye-2026-seeing} \\
TruthfulRAG (knowledge-graph-based factual conflict resolution) & RAG & reasoning & safety & incidental & \cite{liu-2026-truthfulrag} \\
UniKnow / LM-UniKnow & RAG & none & safety & incidental & \cite{kim-2025-uniknow} \\
\emph{Multi-Route Search and Adaptive Fusion (Power QA SLM-guided RAG)} & RAG & none & safety & incidental & \cite{shen-2026-multiroute} \\
\emph{pre-research and post-research verification defenses (verification-enhanced prompting + search-enabled verifier agent) for Deep Research agents} & RAG & none & safety & adversarial & \cite{zhu-2026-deep} \\
XoT (explanation-based thinking framework) & RAG & reasoning & safety & incidental & \cite{zhao-2026-exploring} \\
AgentSecBench Combined defense stack (filtering + projection + redaction + delimiter annotation) & RAG & reasoning & security & adversarial & \cite{alpay-2026-agentsecbench} \\
AttnTrace & RAG & reasoning & security & adversarial & \cite{wang-2025-attntrace} \\
BiRD (Bidirectional Ranking Defense) & RAG & ingestion & security & adversarial & \cite{gao-2026-bird} \\
CachePrune & RAG & reasoning & security & adversarial & \cite{wang-2025-cacheprune} \\
chain-of-defensive-thought (CoDT) & RAG & reasoning & security & adversarial & \cite{wang-2025-chainofdefensivethought} \\
ControlNet (ControlNET) & RAG & ingestion & security & adversarial & \cite{yao-2025-controlnet} \\
FilterRAG / ML-FilterRAG & RAG & ingestion & security & adversarial & \cite{edemacu-2025-defending} \\
InstructDetector & RAG & ingestion & security & adversarial & \cite{wen-2025-defending} \\
LeakSealer & RAG & ingestion & security & adversarial & \cite{panebianco-2025-leaksealer} \\
LLM Firewall (Validator Agent) & RAG & ingestion & security & adversarial & \cite{podpora-2025-llm} \\
MedRAGShield & RAG & ingestion & security & adversarial & \cite{shanid-2026-medragshield} \\
PARSE (Provenance-Aware Retrieval Sanitization) & RAG & ingestion & security & adversarial & \cite{pai-2026-parse} \\
RAG-H & RAG & ingestion & security & adversarial & \cite{ceraku-2026-ragh} \\
ReliabilityRAG & RAG & ingestion & security & adversarial & \cite{shen-2025-reliabilityrag} \\
Rennervate & RAG & reasoning & security & adversarial & \cite{zhong-2025-attention} \\
RobustRAG & RAG & ingestion & security & adversarial & \cite{xiang-2024-certifiably} \\
Salience Normalization (SN-1) & RAG & none & security & adversarial & \cite{zhou-2026-salience} \\
SD-RAG (Selective Disclosure RAG) & RAG & ingestion & security & adversarial & \cite{masoud-2026-sdrag} \\
SecureRag & RAG & ingestion & security & adversarial & \cite{mishra-2025-securerag} \\
TaskTracker (activation-delta task-drift detection) & RAG & reasoning & security & adversarial & \cite{abdelnabi-2024-get} \\
TracLLM & RAG & reasoning & security & adversarial & \cite{wang-2025-tracllm} \\
TriShieldRAG & RAG & ingestion & security & adversarial & \cite{mohanty-2026-trishieldrag} \\
UGCBuster & RAG & ingestion & security & adversarial & \cite{kaya-2025-meets} \\
\emph{policy-first Hybrid RAG architecture (pre-reranking access control + BM25/dense fusion via RRF + evidence-pack citation-grounded generation)} & RAG & none & security & adversarial & \cite{gudipuri-2026-enterprise} \\
\emph{Randomized/Adversarial Suffix Augmentation Training for task-drift probes} & RAG & none & security & adversarial & \cite{rahman-2026-bypassing} \\
\emph{security-aware RAG framework (evidence-based abstention + prompt-injection defense + multi-source agreement verification + trust-weighted retrieval)} & RAG & reasoning & security & both & \cite{kolavennu-2026-securityaware} \\
\emph{Three-layer defense pipeline for RAG PI (embedding content filtering + hierarchical guardrails + response verification)} & RAG & none & security & adversarial & \cite{ramakrishnan-2025-securing} \\
\emph{Three-layer middleware defense for RAG chatbots (Input Screening + Privilege-Constrained Context Assembly + Output Auditing)} & RAG & ingestion & security & adversarial & \cite{saleem-2026-layered} \\
\emph{Three-tier compositional runtime-verification fabric (MonPoly/RTLola-based) for LLM-assisted ISR swarms} & RAG & execution & security & adversarial & \cite{kekatos-2026-missionlevel} \\
\emph{two-plane OPA-based governance architecture for secure RAG (provenance + ABAC + policy-as-code across ingestion/indexing/retrieval/generation)} & RAG & ingestion & security & adversarial & \cite{boggarapu-2026-secure} \\
+ModTool (memory correction tool for MemEvoBench) & memory & none & safety & both & \cite{xie-2026-memevobench} \\
ActiveMem & memory & reasoning & safety & incidental & \cite{jiang-2026-activemem} \\
ADEMA (checkpoint-resumable persistence + segment-level memory condensation + epistemic bookkeeping) & memory & reasoning & safety & incidental & \cite{zhou-2026-adema} \\
Aeon & memory & ingestion & safety & incidental & \cite{arslan-2026-aeon} \\
AFM (Adaptive Focus Memory) & memory & reasoning & safety & incidental & \cite{cruz-2025-adaptive} \\
AgeMem (Agentic Memory) & memory & reasoning & safety & incidental & \cite{yu-2026-agentic} \\
Agent-level affective subsystem (VAD state dynamics) & memory & reasoning & safety & incidental & \cite{subaharan-2026-controlling} \\
AgentFold (folding - granular condensation / deep consolidation context management) & memory & ingestion & safety & incidental & \cite{ye-2025-agentfold} \\
AgentProg (program-guided context management with Execution Tree and global belief state) & memory & ingestion & safety & incidental & \cite{tian-2025-agentprog} \\
AMA-Agent (causality-graph construction + tool-augmented retrieval) & memory & none & safety & incidental & \cite{zhao-2026-amabench} \\
ARC (Active and Reflection-driven Context management) & memory & reasoning & safety & incidental & \cite{yao-2026-arc} \\
AriadneMem & memory & reasoning & safety & incidental & \cite{zhu-2026-ariadnemem} \\
AttriMem & memory & reasoning & safety & incidental & \cite{li-2026-attrimem} \\
AutoMEM & memory & none & safety & incidental & \cite{chen-2026-exploring} \\
Autopilot (gated finite-state-machine execution firewall) & memory & reasoning & safety & incidental & \cite{deng-2026-goalautopilot} \\
CMI (Causal Memory Intervention) & memory & reasoning & safety & both & \cite{srivastava-2026-causal} \\
CoCortex & memory & reasoning & safety & incidental & \cite{kaviyaraj-2026-cocortex} \\
Cognitive Scaffold & memory & reasoning & safety & incidental & \cite{ai-2026-cognitive} \\
CoM (Chain-of-Memory) - Dynamic Memory Chain Evolution + Adaptive Truncation & memory & reasoning & safety & incidental & \cite{xu-2026-chainofmemory} \\
CoMIC & memory & none & safety &  & \cite{wang-2026-comic} \\
Constraint Pinning & memory & ingestion & safety & both & \cite{chen-2026-governance} \\
Context-Agent (dynamic discourse tree) & memory & reasoning & safety & incidental & \cite{hu-2026-contextagent} \\
CWL (Context Window Lifecycle) & memory & ingestion & safety & incidental & \cite{semenov-2026-compaction} \\
D-Mem (Dual-Process Memory System, Quality Gating) & memory & reasoning & safety & incidental & \cite{you-2026-dmem} \\
Dual-Process Memory Architecture (episodic-semantic memory for scientific agents) & memory & reasoning & safety & incidental & \cite{milosevic-2026-episodicsemantic} \\
EMem / EMem-G (event-centric elementary discourse unit memory) & memory & none & safety & incidental & \cite{zhou-2025-simple} \\
Engram (Archive + Research Digest) & memory & none & safety & incidental & \cite{karimi-2026-improving} \\
ENGRAM (typed memory + dense retrieval) & memory & none & safety & incidental & \cite{patel-2025-engram} \\
Ensemble QSP (hierarchical memory architecture) & memory & reasoning & safety & incidental & \cite{tewari-2026-hierarchical} \\
EverMemOS & memory & reasoning & safety & incidental & \cite{hu-2026-evermemos} \\
FGHM (Failure-Gated Hierarchical Memory) with Failure Critic and LTM Pollution Rate (LPR) metric & memory & ingestion & safety & incidental & \cite{liu-2026-failuregated} \\
Fine-Mem (Chunk-level Step Reward + Evidence-Anchored Reward Attribution) & memory & reasoning & safety & incidental & \cite{ma-2026-finemem} \\
FlashMem & memory & none & safety & incidental & \cite{hou-2026-flashmem} \\
Focus (slime-mold-inspired active context compression) & memory & ingestion & safety & incidental & \cite{verma-2026-active} \\
FoldAct & memory & ingestion & safety & incidental & \cite{shao-2025-foldact} \\
GCC (Git-Context-Controller) & memory & ingestion & safety & incidental & \cite{wu-2025-git} \\
GenericAgent (GA) - contextual information density maximization & memory & ingestion & safety & incidental & \cite{liang-2026-genericagent} \\
GRACE (Graph-Regularized Agentic Context Evolution) & memory & reasoning & safety & incidental & \cite{hsu-2026-scoped} \\
GUI-PRA (Process Reward Agent) & memory & none & safety & incidental & \cite{xiong-2025-guipra} \\
HAGE (RL-driven weighted graph evolution) & memory & reasoning & safety & incidental & \cite{jiang-2026-hage} \\
HiAgent & memory & reasoning & safety & incidental & \cite{hu-2024-hiagent} \\
Hierarchical Cognitive Caching (HCC), in ML-Master 2.0 & memory & none & safety & incidental & \cite{zhu-2026-ultralonghorizon} \\
HiMAP-Travel (day-level context isolation + transactional monitor + bargaining protocol) & memory & none & safety & incidental & \cite{bui-2026-himaptravel} \\
HIPIF (Hierarchical Planning and Information Folding) & memory & ingestion & safety & incidental & \cite{diao-2026-hipif} \\
Hippocampus (Dynamic Wavelet Matrix memory) & memory & none & safety & incidental & \cite{li-2026-hippocampus} \\
HMT (Hierarchical Memory Tree) & memory & reasoning & safety & incidental & \cite{tan-2026-enhancing} \\
HORMA (Hierarchical Organize-and-Retrieve Memory Agent) & memory & reasoning & safety & incidental & \cite{hsu-2026-organize} \\
ID-RAG (Identity Retrieval-Augmented Generation) & memory & reasoning & safety & incidental & \cite{platnick-2025-idrag} \\
InftyThink with Cross-Chain Memory (semantic lemma retrieval) & memory & none & safety & incidental & \cite{tekin-2025-directional} \\
Laser (structured action protocol + context register) & memory & ingestion & safety & incidental & \cite{wang-2025-laser} \\
LightMem & memory & none & safety & incidental & \cite{zhang-2026-lightweight} \\
LRE (Learned Relevance Eviction) & memory & none & safety & incidental & \cite{lia-2026-learning} \\
LWM-Planner (Fact Extraction and Compression / FEC) & memory & reasoning & safety & incidental & \cite{holt-2025-factaugmented} \\
M2 (Dynamic Trajectory Summarization + Insight Retrieval Augmentation) & memory & reasoning & safety & incidental & \cite{yan-2026-dualmemory} \\
MAGE (Memory as Agent-Guided Exploration) - hierarchical state tree with Grow/Compress/Maintain/Revise & memory & none & safety & incidental & \cite{chen-2026-semantic} \\
MAGMA (Multi-Graph Agentic Memory Architecture) & memory & reasoning & safety & incidental & \cite{jiang-2026-magma} \\
Mem0\^{}R (unified conversational + repository memory) & memory & none & safety & incidental & \cite{liu-2026-scalable} \\
MEM1 & memory & reasoning & safety & incidental & \cite{zhou-2025-mem1} \\
MemAct (Memory-as-Action) with Dynamic Context Policy Optimization & memory & reasoning & safety & incidental & \cite{zhang-2025-memory} \\
Memanto (typed semantic memory + automated conflict resolution + temporal versioning) & memory & reasoning & safety & incidental & \cite{abtahi-2026-memanto} \\
Memex / MemexRL (indexed experience memory) & memory & reasoning & safety & incidental & \cite{wang-2026-memex} \\
MemFlow (Router Agent + tiered Memory Agent + Validator Agent) & memory & reasoning & safety & incidental & \cite{chen-2026-memflow} \\
MemGUI-Agent / ConAct (Context-as-Action) & memory & none & safety & incidental & \cite{liu-2026-memguiagent} \\
MemMachine (ground-truth-preserving memory) & memory & none & safety & incidental & \cite{wang-2026-memmachine} \\
MemoBrain & memory & reasoning & safety & incidental & \cite{qian-2026-memobrain} \\
Memory Sandbox (user-controlled memory objects) & memory & reasoning & safety &  & \cite{huang-2023-memory} \\
Memory-R1 & memory & reasoning & safety & incidental & \cite{yan-2025-memoryr1} \\
Memory-R2 (LoGo-GRPO training) with shared-parameter fact extractor / memory manager & memory & reasoning & safety & incidental & \cite{yan-2026-memoryr2} \\
MemTrace (memory evolution graph tracing + closed-loop correction) & memory & reasoning & safety & incidental & \cite{deng-2026-memtrace} \\
MemWeaver (graph memory + experience memory + passage memory) & memory & reasoning & safety & incidental & \cite{ye-2026-memweaver} \\
MetaMem (self-evolving meta-memory) & memory & reasoning & safety & incidental & \cite{xin-2026-metamem} \\
MGA (Memory-Driven GUI Agent) - Observer + Structured Memory & memory & reasoning & safety & incidental & \cite{cheng-2025-mga} \\
MMPO (Metacognitive Memory Policy Optimization) with Belief Entropy & memory & reasoning & safety & incidental & \cite{liu-2026-metacognitive} \\
MSCE (Memory-Skill Co-Evolution) & memory & reasoning & safety & incidental & \cite{tang-2026-memory} \\
Multi2 (hierarchical multi-agent decision-making) & memory & none & safety & incidental & \cite{park-2026-multi} \\
Neural Paging / Page Controller & memory & ingestion & safety & incidental & \cite{chen-2026-neural} \\
Oblivion (decay-driven memory control) & memory & reasoning & safety & incidental & \cite{rana-2026-oblivion} \\
OCR-Memory (Optical Context Retrieval Memory) & memory & reasoning & safety & incidental & \cite{li-2026-ocrmemory} \\
PACEvolve - Hierarchical Context Management (HCM) component specifically & memory & none & safety & incidental & \cite{yan-2026-pacevolve} \\
Parallel Compaction & memory & ingestion & safety & incidental & \cite{cim-2026-parallel} \\
PlugMem (knowledge-centric memory graph) & memory & none & safety & incidental & \cite{yang-2026-plugmem} \\
Portable Agent Memory & memory & ingestion & safety &  & \cite{ravindran-2026-portable} \\
PraMem (Practice-derived Experiential Memory) & memory & none & safety & incidental & \cite{li-2026-pramem} \\
PRISM (Pareto-Efficient Retrieval over Intent-Aware Structured Memory) & memory & reasoning & safety & incidental & \cite{peng-2026-prism} \\
PRO-LONG (programmatic memory framework) & memory & reasoning & safety & incidental & \cite{fox-2026-prolong} \\
RAG-Modulo & memory & none & safety & incidental & \cite{jain-2024-ragmodulo} \\
ReAcTree & memory & reasoning & safety &  & \cite{choi-2025-reactree} \\
ReCAP (Recursive Context-Aware Reasoning and Planning) & memory & reasoning & safety & incidental & \cite{zhang-2025-recap} \\
Red-MIRROR & memory & none & safety &  & \cite{khang-2026-redmirror} \\
ReMEmbR (Retrieval-augmented Memory for Embodied Robots) & memory & none & safety & incidental & \cite{anwar-2024-remembr} \\
Rhea (Role-aware Heuristic Episodic Attention) & memory & reasoning & safety & incidental & \cite{hong-2025-rhea} \\
SAM (State-Adaptive Memory) & memory & reasoning & safety & incidental & \cite{hu-2026-sam} \\
Sleep paradigm (Memory Consolidation + Dreaming) & memory & reasoning & safety & incidental & \cite{behrouz-2026-language} \\
SSRP (Self-Synthesizing Reasoning Protocols) & memory & reasoning & safety & incidental & \cite{shehata-2026-attention} \\
STITCH (Structured Intent Tracking in Contextual History) & memory & reasoning & safety & incidental & \cite{yang-2026-grounding} \\
StructAgent (unified causal state + verifier-backed state transitions) & memory & none & safety & incidental & \cite{wu-2026-structagent} \\
Structurally Aligned Subtask-Level Memory & memory & reasoning & safety & incidental & \cite{shen-2026-structurally} \\
SUPO (summarization-based context management for RL) & memory & ingestion & safety & incidental & \cite{lu-2025-scaling} \\
TACT (Think-Act Calibration via activation Steering) & memory & reasoning & safety & incidental & \cite{sui-2026-tact} \\
TDP (Task-Decoupled Planning) & memory & reasoning & safety & incidental & \cite{li-2026-entangled} \\
TokenMizer (typed knowledge-graph session memory) & memory & ingestion & safety & incidental & \cite{mishra-2026-tokenmizer} \\
TokenPilot (Ingestion-Aware Compaction + Lifecycle-Aware Eviction) & memory & ingestion & safety & incidental & \cite{xu-2026-tokenpilot} \\
ToolSelf & memory & none & safety &  & \cite{zhou-2026-toolself} \\
TSUBASA (evolving memory + context distillation) & memory & reasoning & safety & incidental & \cite{zhang-2026-tsubasa} \\
UI-Copilot (memory decoupling + Tool-Integrated Policy Optimization / TIPO) & memory & none & safety & incidental & \cite{lu-2026-uicopilot} \\
UMA (Unified Memory Agent) & memory & reasoning & safety & incidental & \cite{zhang-2026-learning} \\
\emph{Tree-based + hierarchical storage combined memory method (Memory in LLM Era unified framework)} & memory & none & safety & incidental & \cite{wu-2026-memory} \\
WM-SAR (World-Model Subgraph Amplification Repair) & memory & reasoning & safety & incidental & \cite{song-2026-repair} \\
Write-gate sanitization (pre-summarization redaction/rewrite) for memory laundering & memory & ingestion & safety & incidental & \cite{wang-2026-state} \\
Yunque DeepResearch - Dynamic Context Management module specifically & memory & none & safety & incidental & \cite{cai-2026-yunque} \\
Agent Capability Registry & memory & none & security & adversarial & \cite{cai-2026-child} \\
EquiMem & memory & execution & security & adversarial & \cite{meng-2026-equimem} \\
MemAudit (counterfactual memory influence score + memory consistency graph) & memory & execution & security & adversarial & \cite{tan-2026-memaudit} \\
MemLineage & memory & ingestion & security & adversarial & \cite{ouyang-2026-memlineage} \\
ContextRL & cross-modal & reasoning & safety & incidental & \cite{xu-2026-contextaware} \\
Dynamic Contrastive Decoding (DCD) for cross-modality knowledge conflict & cross-modal & none & safety & incidental & \cite{zhu-2024-unraveling} \\
MAGIC-Video (multimodal memory graph + narrative chain) & cross-modal & reasoning & safety & incidental & \cite{li-2026-bridging} \\
ToolScope (Perceive tool within Global Navigator / Agentic Executor / Response Synthesizer) & cross-modal & none & safety & incidental & \cite{deng-2025-toolscope} \\
\emph{DPO-based faithfulness-improvement training for GUI agents against environmental distraction} & cross-modal & none & safety & incidental & \cite{ma-2024-caution} \\
VideoSEAL (decoupled planner-inspector) & cross-modal & reasoning & safety & incidental & \cite{qiu-2026-videoseal} \\
AegisAgent & cross-modal & reasoning & security & adversarial & \cite{wang-2025-aegisagent} \\
ARGUS & cross-modal & reasoning & security &  & \cite{lu-2025-argus} \\
Aura & cross-modal & execution & security & adversarial & \cite{zou-2026-blind} \\
Cross-Agent Multimodal Provenance-Aware Defense Framework & cross-modal & ingestion & security & adversarial & \cite{syed-2025-trustworthy} \\
DE-FIVE & cross-modal & ingestion & security & adversarial & \cite{zhong-2026-defive} \\
Multi-Layer Defense Framework (preprocessing + statistical + neural detection + behavioral monitoring) & cross-modal & ingestion & security & adversarial & \cite{pathade-2025-invisible} \\
PhantomLint & cross-modal & ingestion & security & adversarial & \cite{murray-2025-phantomlint} \\
SmoothVLM & cross-modal & none & security & adversarial & \cite{sun-2024-safeguarding} \\
SnapGuard (visual stability indicator + action-oriented textual signal recovery) & cross-modal & ingestion & security & adversarial & \cite{du-2026-snapguard} \\
SPML-IR specification-based image defense (two-stage input validation + prompt injection detection) & cross-modal & ingestion & security & adversarial & \cite{sharma-2024-defending} \\
\emph{Cross-Agent Multimodal Provenance-Aware Defense Framework (text/visual sanitizer agents + output validator + provenance ledger)} & cross-modal & ingestion & security & adversarial & \cite{syed-2025-crossagent} \\
RS-Claw (hierarchical skill trees for remote-sensing agents) & tool-metadata & none & safety & incidental & \cite{liu-2026-rsclaw} \\
SING (Synthetic Intention Graph) & tool-metadata & none & safety & incidental & \cite{xiao-2026-sing} \\
Skills Architecture (vs. MCP) & tool-metadata & ingestion & safety & incidental & \cite{ayyagari-2026-context} \\
Z-Space & tool-metadata & none & safety &  & \cite{he-2025-zspace} \\
CASCADE & tool-metadata & ingestion & security & adversarial & \cite{turgut-2026-cascade} \\
ContractGuard & tool-metadata & execution & security & adversarial & \cite{iyer-2026-gate} \\
MCP-Guard & tool-metadata & ingestion & security & adversarial & \cite{xing-2025-mcpguard} \\
MCP-Scanner & tool-metadata & ingestion & security & adversarial & \cite{abadeh-2026-mcpscanner} \\
SMCP (Secure Model Context Protocol) & tool-metadata & ingestion & security & adversarial & \cite{hou-2026-smcp} \\
Tool-Guard & tool-metadata & reasoning & security & adversarial & \cite{shi-2026-think} \\
\emph{layered MCP descriptor defense (descriptor integrity verification + pre-context semantic vetting + runtime guardrails)} & tool-metadata & none & security & adversarial & \cite{jamshidi-2025-semantic} \\
ADE-PRF (ADE Predictive Reliability Framework) & multi-agent & none & safety & incidental & \cite{liu-2026-agent} \\
Gated Coordination (Partitioned Information Architecture + Cost-Sensitive Gated Escalation) & multi-agent & none & safety & incidental & \cite{jian-2026-gated} \\
HMAS-2 (Hybrid Multi-Agent System with feedback re-planning) & multi-agent & none & safety & incidental & \cite{chen-2023-scalable} \\
KG-CFR (Knowledge-Grounded Counterfactual Reasoning) & multi-agent & none & safety & incidental & \cite{maslowski-2026-decoupling} \\
LongAgent & multi-agent & reasoning & safety & incidental & \cite{zhao-2024-longagent} \\
MEMO (Memory-augmented MOdel context optimization) & multi-agent & reasoning & safety & incidental & \cite{xie-2026-memo} \\
Meta-Team (collaborative self-evolution) & multi-agent & none & safety & incidental & \cite{hao-2026-evolve} \\
ROMA (Recursive Open Meta-Agents) with GEPA+ & multi-agent & none & safety & incidental & \cite{alzubi-2026-roma} \\
SearchSwarm (delegation-intelligence harness/training for main-agent-to-subagent task decomposition) & multi-agent & ingestion & safety & incidental & \cite{ning-2026-searchswarm} \\
StackPlanner & multi-agent & none & safety & incidental & \cite{zhang-2026-stackplanner} \\
STAR-PólyaMath & multi-agent & none & safety & incidental & \cite{wu-2026-starp} \\
TEP (Textual Equilibrium Propagation) & multi-agent & reasoning & safety & incidental & \cite{chen-2026-textual} \\
Tree of Agents (TOA) & multi-agent & reasoning & safety & incidental & \cite{yu-2025-tree} \\
Agentic JWT (A-JWT) & multi-agent & execution & security & adversarial & \cite{goswami-2025-agentic} \\
BlindGuard & multi-agent & execution & security & adversarial & \cite{miao-2025-blindguard} \\
ChannelGuard & multi-agent & execution & security & adversarial & \cite{hossain-2026-channelguard} \\
Forensic LLM-Trace & multi-agent & none & security & adversarial & \cite{s-2026-forensic} \\
G-Safeguard & multi-agent & execution & security & adversarial & \cite{wang-2025-gsafeguard} \\
GT-MCP (Game-Theoretic Multi-Agent Control) & multi-agent & reasoning & security & adversarial & \cite{jamshidi-2026-gametheoretic} \\
Hybrid monitor scaffold (Monitor Red Teaming / MRT framework) & multi-agent & execution & security & adversarial & \cite{kale-2025-reliable} \\
Maris & multi-agent & execution & security & adversarial & \cite{cui-2025-maris} \\
MAScope & multi-agent & execution & security & adversarial & \cite{wei-2026-input} \\
SafeSieve & multi-agent & ingestion & security & both & \cite{zhang-2025-safesieve} \\
SecureGov-Agent & multi-agent & execution & security & adversarial & \cite{chen-2025-securegovagent} \\
SentinelAgent & multi-agent & execution & security & adversarial & \cite{he-2025-sentinelagent} \\
Token-Level Round-Robin (RR) Collaboration / RRMaj & multi-agent & reasoning & security & adversarial & \cite{liu-2026-consensus} \\
\emph{off-the-shelf LLM behavioral overseer (non-verbal interaction-trace adversarial-intent detection)} & multi-agent & reasoning & security & adversarial & \cite{musaffar-2026-detection} \\
\emph{You've Changed: Detecting Modification of Black-Box Large La} & supply-chain & execution & safety &  & \cite{dima-2025-you} \\
ShieldNet & supply-chain & ingestion & security & adversarial & \cite{yuan-2026-shieldnet} \\
sign-c (server-side signing of execution-bearing response fields) & supply-chain & execution & security & adversarial & \cite{luo-2026-rewriting} \\
Unified Security Orchestration Framework (USOF) & supply-chain & ingestion & security & adversarial & \cite{muthu-2026-unified} \\
SafeClaw-R & skill & execution & safety & both & \cite{wang-2026-safeclawr} \\
\emph{detailed external checklist restatement (vs generic self-check) for long-context skill-following} & skill & none & safety & incidental & \cite{xue-2026-agent} \\
DASGuard & skill & execution & security & adversarial & \cite{tan-2026-prompt} \\
OpenClaw-SRT (Semantic Routing \& Trust) & skill & execution & security & adversarial & \cite{chen-2026-openclawsrt} \\
RouteGuard & skill & ingestion & security & adversarial & \cite{xiao-2026-routeguard} \\
SkillScan & skill & none & security & adversarial & \cite{liu-2026-agentb} \\
SkillSec-Eval defenses (hybrid admission validation, retrieval validation layer, metadata sanitizer, LifecycleGuard taint-tracking, continuous re-validation) & skill & none & security & adversarial & \cite{badhe-2026-agent} \\
SkillVetBench (LLM-as-Judge) & skill & none & security & adversarial & \cite{hossain-2026-skillvetbench} \\
\end{longtable}
\endgroup

